\chapter{Periodic Geometry: Intrinsic Cells and Embedded Planes}
\index{periodic geometry|textbf}
\index{lattice!embedded plane|textbf}
\index{reduced coordinates|textbf}

\label{ch:training-periodic-geometry}

\chapterstatus{pedagogical synthesis and verified software-capability description.
The derivations are illustrative mathematics. The maintained implementation has
software-verification evidence, but this chapter reports no new calculated material
result, numerical convergence claim, scientific validation, or uncertainty
quantification.}

\section{Learning goals}

After working through this chapter, a reader should be able to

\begin{enumerate}
  \item distinguish fractional, intrinsic Cartesian, and ambient Cartesian
        coordinates;
  \item construct the Gram metric of a nonorthogonal primitive basis;
  \item derive the dual reciprocal basis for both square and rectangular direct
        bases;
  \item interpret a one-dimensional Brillouin zone as a reciprocal quotient
        and distinguish physical from reduced momentum;
  \item define a Bloch fiber as a fixed-quasimomentum sector and distinguish
        its quasiperiodic and cell-periodic representations;
  \item explain why an intrinsic two-dimensional cell and a two-dimensional
        plane embedded in three dimensions are different records even when they
        have the same metric;
  \item construct the orthogonal projector onto an embedded plane; and
  \item identify which quantities define the geometry and which should instead
        be reconstructed as derived quantities.
\end{enumerate}

The required background is elementary linear algebra: matrix multiplication,
inner products, rank, inverses of positive-definite matrices, and orthogonal
projection. The chapter prepares the finite-difference construction in
Chapter~\ref{ch:training-bloch-finite-differences} and the controlled
periodic example in Appendix~\ref{app:two-dimensional-reduction}.

\section{Why geometry comes before a periodic discretization}

Chapter~\ref{ch:training-quantum-operators} introduced Hamiltonians abstractly.
Before a periodic derivative can be discretized, a reciprocal vector can be
named, or two represented operators can be compared, we must additionally know
which coordinates and geometry those objects use. A matrix of numbers is not
yet a lattice. It becomes a lattice representation only when we say what its
rows and columns mean and what units its entries carry.

Keep a sheet of ordinary graph paper in mind. On the graph paper, a point can
be located by counting horizontal and vertical grid lines. Those counts play
the role of fractional coordinates. Now imagine stretching the sheet more in
one direction than the other and shearing it so that the grid lines meet at an
oblique angle. The counting coordinates did not change: one still moves by
``one cell'' when one count increases by one. Physical distances did change.
The matrix $A$ records that change from counting coordinates to Cartesian
vectors.

Now lift the same sheet, without bending it, and rotate it in a three-dimensional
room. Distances measured on the sheet do not change, but the sheet has acquired
an orientation relative to the room. This is the distinction between intrinsic
geometry and embedding data. The Gram metric records the distances and angles
available to an observer confined to the sheet. The ambient basis, frame, and
projector additionally record how the sheet sits in the room.

That picture is useful because the software boundary follows the same logic.
An intrinsic scalar differential operator needs the metric. A later model that
couples the plane to a laboratory-axis electric field also needs the ambient
orientation. If we discard the orientation when constructing the intrinsic
operator, we may not be able to reconstruct the coupling correctly. If we pass
all ambient data into every scalar stencil routine, on the other hand, we make
a simple intrinsic algorithm depend on information it does not mathematically
need. The adapter introduced later in this lecture preserves both sides of the
boundary.

A recurring discipline of these notes is therefore
\begin{equation}
  \text{coordinates}
  \longrightarrow
  \text{geometry}
  \longrightarrow
  \text{represented operator}
  \longrightarrow
  \text{physical model}.
  \label{eq:training-geometry-layer-sequence}
\end{equation}
The arrow is not permission to collapse the objects on either side of it. Each
step adds meaning and assumptions.

\section{Primitive vectors as a coordinate map}
\index{primitive vector}
\index{fractional coordinate}

Let a $d$-dimensional periodic cell have ordered primitive vectors
$\mathbf a_1,\ldots,\mathbf a_d$. Store those vectors as the columns of a
matrix
\begin{equation}
  A=
  \begin{bmatrix}
    \mathbf a_1 & \cdots & \mathbf a_d
  \end{bmatrix}.
  \label{eq:training-direct-basis}
\end{equation}
A fractional coordinate $\mathbf s\in\mathbb R^d$ maps to Cartesian position
through
\begin{equation}
  \mathbf r=A\mathbf s.
  \label{eq:training-fractional-to-cartesian}
\end{equation}
The entries of $\mathbf s$ are dimensionless. The columns of $A$ carry the
length unit. Consequently, $\mathbf r$ carries that same length unit.

Two cases must be distinguished.

\begin{description}
  \item[Intrinsic two-dimensional basis.] $A\in\mathbb R^{2\times2}$ maps two
    fractional coordinates into an intrinsic Cartesian plane.
  \item[Two-dimensional plane embedded in three dimensions.]
    $A\in\mathbb R^{3\times2}$ maps the same number of fractional coordinates
    into ambient three-space.
\end{description}

Both cases describe two periodic directions because both matrices have two
columns. The number of rows instead records the ambient Cartesian dimension.
A $3\times2$ basis therefore does not define a third periodic direction.

\subsection{Three coordinate languages}

During a calculation we will move repeatedly among three related coordinate
languages. It is worth naming them before doing any algebra.

\begin{description}
  \item[Fractional coordinates $\mathbf s$.]
    These say how many primitive vectors are used. Their components are
    dimensionless, and integer translations identify equivalent points in a
    periodic cell.
  \item[Intrinsic Cartesian coordinates $\mathbf x$.]
    These are ordinary orthonormal coordinates inside the plane. They carry a
    length unit. Choosing them requires an in-plane orthonormal frame.
  \item[Ambient Cartesian coordinates $\mathbf r$.]
    These are coordinates in the surrounding space. For an embedded plane,
    $\mathbf r$ has three components even though only two independent
    coordinates are periodic.
\end{description}

The same physical in-plane point can consequently be written in several ways:
\begin{equation}
  \mathbf r=A\mathbf s=Q\mathbf x,
  \qquad
  \mathbf x=Q^{\mathsf T}A\mathbf s.
  \label{eq:training-three-coordinate-languages}
\end{equation}
Nothing physical changes when we translate the point faithfully between these
languages. What changes is the representation and therefore the matrices that
act on coordinate arrays.

A good notebook habit is to annotate every vector at first use. For example,
write ``$\mathbf s$: fractional, dimensionless, two components'' and
``$\mathbf r$: ambient Cartesian, angstrom, three components.'' Such notes may
look excessive in a two-line exercise. They become invaluable when a retained
artifact is read months later or when a two-dimensional model is passed through
a three-dimensional file interface.

\begin{remark}[Columns, not rows]
The convention in Equation~\eqref{eq:training-direct-basis} is part of the
representation. Transposing $A$ without changing the formulas changes the
coordinate map. A retained record or file should state the vector-axis
convention rather than relying on visual inspection of a rectangular array.
\end{remark}

\section{The induced metric}
\index{Gram metric|textbf}
\index{metric tensor!direct lattice}

The Cartesian squared length of a displacement is
\begin{align}
  \lVert \mathrm{d}\mathbf r\rVert_2^2
  &= (A\,\mathrm{d}\mathbf s)^{\mathsf T}
     (A\,\mathrm{d}\mathbf s) \\
  &= \mathrm{d}\mathbf s^{\mathsf T}G\,\mathrm{d}\mathbf s,
  \label{eq:training-induced-metric}
\end{align}
where
\begin{equation}
  G=A^{\mathsf T}A.
  \label{eq:training-gram-metric}
\end{equation}
The matrix $G$ is the Gram metric of the ordered primitive vectors. It contains
all intrinsic lengths and angles:
\begin{equation}
  G_{ij}=\mathbf a_i\cdot\mathbf a_j.
  \label{eq:training-metric-components}
\end{equation}
If the columns of $A$ are linearly independent, then $G$ is symmetric positive
definite and $G^{-1}$ exists.

This observation explains why an intrinsic $2\times2$ basis and an embedded
$3\times2$ basis can define the same scalar periodic geometry. If
$A_{\mathrm{int}}^{\mathsf T}A_{\mathrm{int}}
=A_{\mathrm{emb}}^{\mathsf T}A_{\mathrm{emb}}$, they assign the same length to
every fractional displacement. The embedded basis nevertheless carries extra
orientation information that the metric alone cannot recover.

\subsection{Orthogonal square example}

For
\begin{equation}
  A=
  \begin{bmatrix}
    a & 0\\
    0 & a\\
    0 & 0
  \end{bmatrix},
  \label{eq:training-square-embedded-basis}
\end{equation}
we obtain
\begin{equation}
  G=a^2I_2,
  \qquad
  G^{-1}=a^{-2}I_2.
  \label{eq:training-square-metric}
\end{equation}
Fractional coordinate lines are orthogonal and have equal physical scale.
This is the geometry used by the original square-lattice controlled model in
Appendix~\ref{app:two-dimensional-reduction}.

\subsection{Skew-cell example}

Let
\begin{equation}
  \mathbf a_1=a
  \begin{bmatrix}1\\0\\0\end{bmatrix},
  \qquad
  \mathbf a_2=a
  \begin{bmatrix}\cos\theta\\\sin\theta\\0\end{bmatrix},
  \qquad 0<\theta<\pi.
  \label{eq:training-skew-basis}
\end{equation}
Then
\begin{equation}
  G=a^2
  \begin{bmatrix}
    1 & \cos\theta\\
    \cos\theta & 1
  \end{bmatrix}.
  \label{eq:training-skew-metric}
\end{equation}
The off-diagonal entries are geometric: they record that the coordinate axes
are not orthogonal. They do not, by themselves, describe an anisotropic
material response.

\subsection{A fully worked sixty-degree cell}

Let us pause over one cell long enough to make the abstractions concrete. Take
\begin{equation}
  \mathbf a_1=a
  \begin{bmatrix}1\\0\\0\end{bmatrix},
  \qquad
  \mathbf a_2=a
  \begin{bmatrix}1/2\\\sqrt{3}/2\\0\end{bmatrix}.
  \label{eq:training-sixty-degree-basis}
\end{equation}
The two primitive vectors have length $a$ and meet at sixty degrees. Their Gram
metric is
\begin{equation}
  G=a^2
  \begin{bmatrix}
    1 & 1/2\\
    1/2 & 1
  \end{bmatrix},
  \qquad
  G^{-1}=\frac{1}{a^2}
  \begin{bmatrix}
    4/3 & -2/3\\
    -2/3 & 4/3
  \end{bmatrix}.
  \label{eq:training-sixty-degree-metric}
\end{equation}
Notice the negative off-diagonal entries of the inverse metric. They do not
mean that a squared length can be negative. Positive definiteness concerns the
whole quadratic form
$\mathrm{d}\mathbf s^{\mathsf T}G\,\mathrm{d}\mathbf s$, not the sign of
each matrix entry.

Consider three fractional displacements:
\begin{equation}
  \begin{aligned}
  \mathrm{d}\mathbf s_1&=(1,0)^{\mathsf T},\\
  \mathrm{d}\mathbf s_2&=(0,1)^{\mathsf T},\\
  \mathrm{d}\mathbf s_3&=(1,-1)^{\mathsf T}.
  \end{aligned}
\end{equation}
Equation~\eqref{eq:training-induced-metric} assigns all three the same physical
length $a$. The third result is a useful check: in a sixty-degree rhombus, the
difference between the two primitive vectors is also an edge of the associated
equilateral triangle.

Now rotate both vectors in ambient three-space. The individual Cartesian
components change, but every entry of $G$ remains the same. An intrinsic scalar
calculation should not change merely because the flat sample was rotated on a
laboratory table. A coupling to a fixed laboratory direction may change, which
is exactly why we retain the ambient representation separately.

The following small program turns the board derivation into explicit arrays.
The imports are shown because a reader should always be able to identify which
package owns each numerical operation. Inline comments explain why a line is
present rather than merely restating its syntax.

\begin{pythonexample}{checking an embedded periodic plane}
import numpy as np

# Store primitive vectors as columns of one ambient 3-by-2 basis.
a = 2.0
A = a * np.array(
    [[1.0, 0.5], [0.0, np.sqrt(3.0) / 2.0], [0.0, 0.0]],
    dtype=np.float64,
)

# Derive intrinsic geometry from the defining basis.
G = A.T @ A
G_inverse = np.linalg.inv(G)
B = 2.0 * np.pi * A @ G_inverse
P = A @ G_inverse @ A.T

# Measure duality and orthogonal-projector identities independently.
duality_error = np.max(np.abs(A.T @ B - 2.0 * np.pi * np.eye(2)))
symmetry_error = np.max(np.abs(P.T - P))
idempotence_error = np.max(np.abs(P @ P - P))

print("Gram metric:\n", G)
print("projector rank:", np.linalg.matrix_rank(P))
print("maximum duality error:", duality_error)
print("maximum projector symmetry error:", symmetry_error)
print("maximum projector idempotence error:", idempotence_error)
\end{pythonexample}

This code is pedagogical synthetic data, not a material calculation.

\documentcallout{Notebook to do}{Create
\path{examples/tutorials/research-monograph/foundations/periodic_geometry.ipynb}
with explicit imports, explanatory inline comments, physical units, transparent
NumPy checks, and exercises using the maintained PhysKit adapter and strict
codec. Keep the notebook unexecuted in the repository and verify it through a
separate bounded test route.}

\section{Dual reciprocal vectors}
\index{reciprocal lattice!dual basis}

For a full-column-rank $A\in\mathbb R^{m\times d}$, with $m\geq d$, define
\begin{equation}
  B=2\pi A\left(A^{\mathsf T}A\right)^{-1}
   =2\pi A G^{-1}.
  \label{eq:training-rectangular-dual}
\end{equation}
Its columns $\mathbf b_j$ obey
\begin{equation}
  A^{\mathsf T}B=2\pi I_d,
  \qquad
  \mathbf a_i\cdot\mathbf b_j=2\pi\delta_{ij}.
  \label{eq:training-dual-identity}
\end{equation}
When $A$ is square and invertible, Equation~\eqref{eq:training-rectangular-dual}
reduces to the familiar expression
\begin{equation}
  B=2\pi A^{-\mathsf T}.
  \label{eq:training-square-dual}
\end{equation}
The rectangular formula is therefore not a different reciprocal-lattice
convention. It is the full-column-rank extension of the same duality condition.

\subsection{Deriving the rectangular formula at the board}

The defining requirement is not ``take an inverse.'' The defining requirement
is the duality relation
\begin{equation}
  A^{\mathsf T}B=2\pi I_d.
\end{equation}
For a rectangular $A$, an ordinary inverse does not exist. We therefore ask for
columns of $B$ that lie in the same plane as the columns of $A$. Any such
candidate has the form $B=AC$ for some $d\times d$ coefficient matrix $C$.
Substitution gives
\begin{equation}
  A^{\mathsf T}AC=GC=2\pi I_d.
\end{equation}
Full column rank makes $G$ invertible, so
\begin{equation}
  C=2\pi G^{-1}
\end{equation}
and hence $B=2\pi AG^{-1}$. This short derivation tells us more than a memorized
formula. It tells us why the dual vectors remain in the tangent plane and why
rank loss makes the construction fail.

For the sixty-degree cell in
Equation~\eqref{eq:training-sixty-degree-basis}, the result is
\begin{equation}
  B=\frac{2\pi}{a}
  \begin{bmatrix}
    1 & 0\\
    -1/\sqrt{3} & 2/\sqrt{3}\\
    0 & 0
  \end{bmatrix}.
  \label{eq:training-sixty-degree-dual}
\end{equation}
A quick multiplication by $A^{\mathsf T}$ recovers $2\pi I_2$. Performing that
multiplication is preferable to trusting the visual appearance of the vectors.

Because $A$ carries length, $G$ carries length squared, $G^{-1}$ carries inverse
length squared, and $B$ carries inverse length. These unit checks often reveal a
transposition or scaling mistake before a numerical calculation is run. In
lecture, it is useful to say the units aloud while moving through the product:
``length times inverse length squared gives inverse length.''

\section{The one-dimensional reciprocal quotient}
\index{Brillouin zone!one-dimensional}
\index{reduced momentum}

The general dual-basis construction becomes especially transparent in one
dimension. Let the direct primitive vector have length $a>0$. Its reciprocal
primitive vector has magnitude
\begin{equation}
  G=\frac{2\pi}{a},
  \label{eq:training-geometry-one-dimensional-reciprocal}
\end{equation}
because $aG=2\pi$. A physical wave vector $k$ can be written in the
dimensionless reduced coordinate
\begin{equation}
  \kappa=\frac{k}{G}.
  \label{eq:training-geometry-reduced-momentum}
\end{equation}
Wave vectors whose reduced coordinates differ by an integer occupy the same
point of the reciprocal quotient. Thus the first Brillouin zone is not
fundamentally a closed interval with two distinct endpoints. It is a circle of
reciprocal classes that can be represented, for example, by
\begin{equation}
  -\frac12\leq\kappa<\frac12.
  \label{eq:training-geometry-first-brillouin-zone}
\end{equation}
The endpoint $1/2$ is excluded because it is equivalent to $-1/2$ after one
reciprocal translation. Including both may be useful when displaying a closed
plotting path, but it duplicates one point when constructing a complete
periodic mesh.

\subsection{Bloch sectors and fiber operators}
\index{Bloch fiber|textbf}
\index{direct integral}
\index{quasimomentum sector}

The reciprocal quotient labels more than plotting coordinates. Suppose that
an operator $\widehat H$ on the infinite periodic system commutes with
translation by $a$. Bloch theory resolves its state space into sectors labelled
by a reciprocal class $[k]$ \parencite{kittel2004,simon2013}. In the sector
labelled by $k$, a wavefunction
obeys the quasiperiodic condition
\begin{equation}
  \psi_k(x+a)=e^{ika}\psi_k(x).
  \label{eq:training-geometry-bloch-condition}
\end{equation}
Changing $k$ by an integer multiple of $G$ gives the same reciprocal class and
therefore does not create a new physical sector.

A \emph{Bloch fiber space} is the state space belonging to one fixed reciprocal
class $[k]$. The corresponding \emph{Bloch-fiber operator} is the action of
$\widehat H$ restricted to that sector. This is an operator and state-space
statement, not yet a choice of numerical basis or finite matrix.

There are two common, unitarily related ways to represent the same fiber. One
may retain the original wavefunction on a primitive cell and impose
Equation~\eqref{eq:training-geometry-bloch-condition}. Alternatively, write
\begin{equation}
  \psi_k(x)=e^{ikx}u_k(x),
  \qquad
  u_k(x+a)=u_k(x),
  \label{eq:training-geometry-bloch-factorization}
\end{equation}
and work in the fixed cell-periodic space
$\mathcal H_{\mathrm{per}}=L^2_{\mathrm{per}}([0,a])$. In this representation,
the fiber operator is
\begin{equation}
  \widehat H(k)
  =e^{-ikx}\widehat H e^{ikx}
  \quad\text{on }\mathcal H_{\mathrm{per}},
  \label{eq:training-geometry-fiber-definition}
\end{equation}
with an operator domain carrying periodic boundary conditions. For example,
if
\begin{equation}
  \widehat H
  =-\frac{\hbar^2}{2m}\frac{\mathrm d^2}{\mathrm d x^2}+V(x),
  \qquad V(x+a)=V(x),
\end{equation}
then
\begin{equation}
  \widehat H(k)
  =\frac{\hbar^2}{2m}
   \left(-i\frac{\mathrm d}{\mathrm d x}+k\right)^2+V(x).
  \label{eq:training-geometry-schrodinger-fiber}
\end{equation}
The $k$ dependence therefore belongs to the represented derivative and its
domain convention; it is not an additional potential.

At the operator level, the Floquet--Bloch transformation expresses this
sector decomposition as a direct integral \parencite{kuchment2016},
\begin{equation}
  \widehat H
  \simeq
  \int_{[-G/2,G/2)}^{\oplus}
    \widehat H(k)\,\frac{a\,\mathrm d k}{2\pi}.
  \label{eq:training-geometry-direct-integral}
\end{equation}
The symbol $\simeq$ denotes unitary equivalence, and $\oplus$ indicates that
there is one fiber space at each reciprocal class. Solving one fiber produces
energies and states at one $k$; sampling many fibers produces a band
representation. A plane-wave cutoff or finite-difference grid then produces a
finite representation of each fiber. Such a matrix approximates or represents
the declared fiber operator under a stated discretization; it is not the
continuum fiber operator itself.

For a particle of mass $m$, the reciprocal vector also defines the energy
scale
\begin{equation}
  E_G=\frac{\hbar^2G^2}{2m}.
  \label{eq:training-geometry-reciprocal-energy}
\end{equation}
An energy written as $E/E_G$ and a potential amplitude written as $V_0/E_G$
are dimensionless only after $a$, $G$, $m$, and the scaling convention have
been declared. Chapter~\ref{ch:training-isolated-band-fourier-reduction} uses
these coordinates and scales to derive a complete one-dimensional
Bloch-to-lattice reduction.

\section{Projection and coordinate recovery}
\index{pseudoinverse}
\index{projection!ambient plane}

The left pseudoinverse of a full-column-rank embedded basis is
\begin{equation}
  A^{+}=G^{-1}A^{\mathsf T}.
  \label{eq:training-left-pseudoinverse}
\end{equation}
For an in-plane point $\mathbf r=A\mathbf s$,
\begin{equation}
  A^{+}\mathbf r=\mathbf s.
  \label{eq:training-fractional-recovery}
\end{equation}
For a general ambient vector, $A^{+}\mathbf r$ returns the fractional
coordinates of its orthogonal in-plane projection. The associated ambient
projector is
\begin{equation}
  P=AA^{+}=AG^{-1}A^{\mathsf T}.
  \label{eq:training-ambient-projector}
\end{equation}
It satisfies
\begin{equation}
  P^{\mathsf T}=P,
  \qquad
  P^2=P,
  \qquad
  \operatorname{Ran}P=\operatorname{Ran}A.
  \label{eq:training-projector-identities}
\end{equation}
An out-of-plane component is removed by this operation. That removal must be
explicit: it is not evidence that the ambient vector was originally confined
to the plane.

\subsection{A projection thought experiment}

Suppose the plane is the ordinary $xy$ plane and an ambient vector is
$\mathbf v=(v_x,v_y,v_z)^{\mathsf T}$. The projector returns
$(v_x,v_y,0)^{\mathsf T}$. No information inside the projected vector tells us
whether $v_z$ was originally zero or was discarded. This is why a projected
result should retain provenance linking it to its ambient source and to the
projection action.

The same issue is easy to miss for a tilted plane because none of the ambient
components need vanish. The equation $P\mathbf v=\mathbf v$ is the appropriate
test for whether $\mathbf v$ already lies in the plane. Merely observing three
nonzero components says nothing either way.

The distinction has a physical counterpart. An ambient electric field may have
both tangent and normal components. A strictly two-dimensional model can use
the tangent component only after a modeling decision has said what happens to
the normal component. Linear algebra can perform the projection, but it cannot
make that physical decision for us.

\section{An intrinsic frame that preserves the embedding}
\index{orthonormal frame!embedded plane}

Let $Q\in\mathbb R^{3\times2}$ contain an ordered orthonormal frame for the
plane, so that
\begin{equation}
  Q^{\mathsf T}Q=I_2,
  \qquad
  QQ^{\mathsf T}=P.
  \label{eq:training-plane-frame}
\end{equation}
The coordinates of the primitive vectors in that frame form the intrinsic
basis
\begin{equation}
  C=Q^{\mathsf T}A\in\mathbb R^{2\times2}.
  \label{eq:training-intrinsic-adapter}
\end{equation}
The intrinsic basis has exactly the same Gram metric:
\begin{equation}
  C^{\mathsf T}C
  =A^{\mathsf T}QQ^{\mathsf T}A
  =A^{\mathsf T}A
  =G.
  \label{eq:training-adapter-metric-identity}
\end{equation}
This is the key adapter identity. Algorithms that depend only on the intrinsic
metric can operate on $C$, while the original $A$, $Q$, and $P$ retain the
ambient orientation needed by later composition.

If the ambient basis is rigidly rotated to $RA$ with
$R^{\mathsf T}R=I_3$, then
\begin{equation}
  (RA)^{\mathsf T}(RA)=G.
  \label{eq:training-rotation-invariance}
\end{equation}
A geometric scalar Laplacian is therefore invariant under ambient rigid
rotation. Coupling to an ambient vector field need not be invariant unless the
field is rotated consistently.

\section{An embedded plane is not a slab}
\index{slab!distinction from embedded plane}

A flat embedded plane and a slab answer different questions.

\begin{description}
  \item[$3\times2$ embedded plane.] Two periodic coordinates; no third direct
    primitive vector; intrinsic differential operators act only in the plane.
  \item[$3\times3$ slab or bulk cell.] Three direct primitive vectors. Even if
    one direction contains vacuum or is sampled sparsely, it remains part of
    the declared three-dimensional cell and boundary-condition model.
\end{description}

This distinction is also separate from the auxiliary third direction used by
some historical Wannier90 interfaces in
Appendix~\ref{app:two-dimensional-reduction}. The phrase ``embedding'' there
can refer to satisfying a three-dimensional interface convention. The present
chapter instead concerns the geometric map from an intrinsically
two-dimensional lattice into ambient Cartesian three-space.

\subsection{A language test}

Whenever the word ``two-dimensional'' appears, ask two questions.

\begin{enumerate}
  \item How many independent coordinates does the state depend on?
  \item In how many Cartesian dimensions is the geometry represented?
\end{enumerate}

For an intrinsic plane the answers are two and two. For an embedded plane they
are two and three. For a three-dimensional slab model they are three and three,
even if the third direction is physically long or contains vacuum. This
language test prevents a data-shape observation from silently becoming a
physical claim.

\section{Maintained computational representation}
\index{PhysKit!embedded lattice}
\index{serialization!embedded lattice}

The maintained PhysKit dependency represents the defining basis with
\path{EmbeddedDirectLattice2D}. An explicit adapter constructs the intrinsic
basis, induced metric, inverse metric, ambient dual reciprocal basis,
orthonormal frame, and tangent-plane projector. The adapter preserves the
original $A$ rather than silently replacing it by $C$.

The corresponding strict JSON codec serializes only the defining $3\times2$
basis, its unit, a schema version, a representation discriminator, and the
column-vector convention. It deliberately does not serialize $G$, $G^{-1}$,
$B$, $Q$, or $P$. Those quantities are recomputed and checked after decoding.
This design prevents a retained document from carrying mutually inconsistent
defining and derived geometry.

The accepted dependency revision for the software-capability description in
this draft is
\path{949d106bc98975a18ad86d6fa84092ef0c68298e}. A consuming-project smoke test
constructs a physical $3\times2$ basis, performs the strict JSON round trip,
adapts it to the intrinsic metric, and constructs a Hermitian two-dimensional
Bloch-periodic scalar Laplacian. That test establishes bounded software
composition only.

\section{How to record a geometry laboratory}
\index{laboratory!periodic geometry}

A reproducible geometry exercise should retain more than the final metric.
Before computing, write down the following items:

\begin{enumerate}
  \item the ordered primitive vectors and whether they are stored as rows or
        columns;
  \item the ambient and intrinsic dimensions;
  \item the length unit of the direct basis;
  \item the fractional-coordinate origin and cell convention;
  \item the intended interpretation: intrinsic plane, embedded plane, slab, or
        bulk cell; and
  \item the invariants that will be checked independently.
\end{enumerate}

After computing, retain at least the defining basis and observations of the
checks
\begin{equation}
  G=A^{\mathsf T}A,
  \qquad
  A^{\mathsf T}B=2\pi I,
  \qquad
  P^{\mathsf T}=P,
  \qquad
  P^2=P.
  \label{eq:training-geometry-lab-checks}
\end{equation}
Derived arrays may be useful in a report, but they should not be allowed to
compete with the defining basis as independent authorities. If a serialized
metric and a serialized basis disagree, the document has not preserved one
coherent geometry.

For a classroom calculation, print a few small matrices and inspect them. For a
maintained research artifact, use typed records, strict decoding, explicit
units, and a manifest. The scale changes; the conceptual obligation does not.

\section{End-of-lecture synthesis}

The central object of this lecture is the map $\mathbf r=A\mathbf s$. From it we
obtained four families of derived information:

\begin{enumerate}
  \item $G=A^{\mathsf T}A$ describes intrinsic lengths and angles;
  \item $B=2\pi AG^{-1}$ supplies reciprocal vectors dual to the direct basis
        and thereby defines reduced reciprocal coordinates;
  \item $A^+=G^{-1}A^{\mathsf T}$ recovers fractional coordinates of the
        in-plane projection; and
  \item $P=AG^{-1}A^{\mathsf T}$ records the plane inside the ambient space.
\end{enumerate}

The formulas are short, but the distinctions are durable. Fractional
coordinates are not Cartesian coordinates. Intrinsic geometry is not ambient
orientation. A flat embedded plane is not a slab. A geometric off-diagonal
metric entry is not a material effective-mass tensor. These distinctions are
the vocabulary needed for the next lecture, where the metric becomes a
differential operator.

\section{Exercises}

\begin{enumerate}
  \item For the skew basis in Equation~\eqref{eq:training-skew-basis}, compute
        $G^{-1}$ and show explicitly that it diverges as
        $\theta\rightarrow0$ or $\theta\rightarrow\pi$. Interpret the result as
        loss of column rank rather than as a material singularity.
  \item Starting from Equation~\eqref{eq:training-rectangular-dual}, prove
        Equation~\eqref{eq:training-dual-identity} without assuming that $A$ is
        square.
  \item Show that the projector in
        Equation~\eqref{eq:training-ambient-projector} is unchanged when $A$ is
        replaced by $AS$ for any invertible $2\times2$ matrix $S$. Explain what
        does and does not remain the same about the periodic cell.
  \item Choose two nonorthogonal vectors in three dimensions, construct $A$,
        $G$, $B$, and $P$, and check the duality and projector identities
        numerically.
  \item Explain why retaining only $G$ is sufficient for a scalar intrinsic
        Laplacian but insufficient for coupling an in-plane model to a declared
        ambient electric-field direction.
  \item Design one malformed JSON case for each of the following failures:
        wrong vector-axis declaration, duplicate field, nonfinite number,
        numeric string, incompatible unit, and rank-deficient basis. State
        which boundary should reject each case.
\end{enumerate}

\section{Evidence boundary}

The identities in this chapter are geometry and linear algebra. Passing the
associated software tests establishes implementation behavior for the tested
contracts. It does not establish convergence of a finite-difference model,
adequacy of a kinetic Hamiltonian, a material lattice, a slab approximation,
or a semiconductor effective-mass tensor. Those questions require the
separate operators, calculations, and evidence developed later in the
monograph.
